% Adding the MATLAB feedback controller to the gantry derivation.
% Notation follows docs/writeup/jan-augmentation-writeup.tex. House style follows
% docs/writeup/jan-writeup-guideline.md. Compile: pdflatex controller-in-derivation.tex
\documentclass[11pt]{article}

\usepackage[a4paper,margin=2.2cm]{geometry}
\usepackage[T1]{fontenc}
\usepackage[utf8]{inputenc}
\usepackage{lmodern}
\usepackage{amsmath,amssymb,mathtools,bm}
\usepackage{microtype}
\usepackage{enumitem}
\usepackage{titlesec}
\usepackage{booktabs}
\usepackage{graphicx}
\usepackage[hidelinks]{hyperref}

\titleformat{\section}{\Large\bfseries}{\thesection}{0.75em}{}
\titleformat{\subsection}{\large\bfseries}{\thesubsection}{0.75em}{}
\titlespacing*{\section}{0pt}{1.4ex plus .2ex minus .2ex}{0.9ex}
\titlespacing*{\subsection}{0pt}{1.0ex plus .2ex minus .2ex}{0.5ex}

\newcommand{\R}{\mathbb{R}}
\newcommand{\diag}{\operatorname{diag}}
\newcommand{\xb}{\tilde{x}}
\newcommand{\yh}{\hat{y}}
\newcommand{\Cfb}{C_{\mathrm{fb}}}
\newcommand{\Gop}{G_{\mathrm{op}}}
\newcommand{\So}{S_{\mathrm{o}}}
\newcommand{\Si}{S_{\mathrm{i}}}
\newcommand{\Cnorm}{C_{\mathrm{n}}}
\newcommand{\fms}{f_{\mathrm{ms}}}
\newcommand{\ufb}{u_{\mathrm{fb}}}

\setlist{nosep, leftmargin=1.4em}

\begin{document}

\begin{center}
  {\LARGE\bfseries Adding the feedback controller to the derivation}\\[0.4em]
  {\large Dirk van den Berg}\\[0.2em]
  {\small Notation follows \texttt{docs/writeup/jan-augmentation-writeup.tex}. Every equation
  here is realised in \texttt{verify\_controller.py} and gated against the generated records.}
\end{center}
\vspace{0.4em}

\section{Purpose and scope}

The trajectory records are generated in closed loop and replayed open loop. This note writes the
controller that closes that loop as an explicit transfer function, so that the loop can be
carried through the derivation as equations rather than as a MATLAB object. Section~\ref{sec:ctrl}
gives the controller, Section~\ref{sec:loop} solves the loop, Section~\ref{sec:conseq} states
what the loop does and does not change about the model comparison, and Section~\ref{sec:gate}
reports the numerical gate.

\medskip\noindent\textbf{Symbols that clash with the baseline.}\quad The baseline write-up already
uses $C$ for the damping matrix of Eq.~(5) there and $K$ for the stiffness matrix. The controller
and its normalisation gain therefore take different symbols in this note:

\begin{center}
\begin{tabular}{lll}
\toprule
symbol & meaning & clashes with \\
\midrule
$\Cfb(z)$ & discrete feedback controller, $3\times3$ diagonal & $C$, damping matrix \\
$\Cnorm(s)$ & the shared normalised controller shape & \\
$\kappa_j$ & per-channel normalisation gain & $K$, stiffness matrix \\
$\omega_{\mathrm{b}}$ & design bandwidth in rad/s & \\
\bottomrule
\end{tabular}
\end{center}

\section{Where the controller attaches}

The measured output of the baseline is the vector of stage positions,
\begin{equation}
  q = [\,X,\Theta,Y\,]^\top , \qquad
  \yh = P^\top q = [\,X_1,X_2,Y\,]^\top , \qquad
  P =
  \begin{bmatrix}
    1 & 1 & 0 \\[1mm]
    \tfrac{L_b}{2} & -\tfrac{L_b}{2} & 0 \\[1mm]
    0 & 0 & 1
  \end{bmatrix} ,
  \label{eq:frame}
\end{equation}
and the controller acts on $\yh$, that is, in the \emph{stage} frame. This fixes the frame for
everything that follows: the loop is diagonal in stage coordinates and the logical frame appears
only inside the plant. The plant seen by the loop is therefore the stage-to-stage map
\begin{equation}
  G(s) : u \mapsto \yh , \qquad
  M(Y)\,\ddot q + C\,\dot q + K\,q = P^\top u , \qquad \yh = P^\top q ,
  \label{eq:plant}
\end{equation}
with $u = [\,F_1,F_2,F_y\,]^\top$ the stage forces of the baseline write-up.

\section{The design plant is frozen}

The controller is designed once per record on a \emph{frozen} operating point $Y_{\mathrm{op}}$
and with the \emph{full} payload mass $m_h$, so it is linear time invariant even though the plant
is not. Write
\begin{equation}
  M_{\mathrm{op}} =
  \begin{bmatrix}
    \alpha                            & \beta - m_h Y_{\mathrm{op}}                & 0 \\[2mm]
    \beta - m_h Y_{\mathrm{op}}       & \gamma + m_h Y_{\mathrm{op}}^2             & -m_h d \\[2mm]
    0                                 & -m_h d                                     & m_h
  \end{bmatrix},
  \label{eq:Mop}
\end{equation}
with $\alpha,\beta,\gamma$ the constants of the baseline write-up. The design plant in state
space, stage in and stage out, is
\begin{equation}
  A_{\mathrm{op}} =
  \begin{bmatrix}
    0_{3\times3} & I_3 \\[2mm]
    -M_{\mathrm{op}}^{-1}K & -M_{\mathrm{op}}^{-1}C
  \end{bmatrix},
  \qquad
  B_{\mathrm{op}} =
  \begin{bmatrix}
    0_{3\times3} \\[2mm]
    M_{\mathrm{op}}^{-1}P
  \end{bmatrix},
  \qquad
  C_{\mathrm{op}} =
  \begin{bmatrix}
    P^\top & 0_{3\times3}
  \end{bmatrix},
  \label{eq:designss}
\end{equation}
and $\Gop(z)$ is \eqref{eq:designss} discretised with a zero-order hold at the sample time
$t_s = 5\cdot10^{-5}$~s. Note that $\Gop$ is the \emph{rigid} baseline: it carries neither the
absorber of the truth nor the payload split, so it is a design model only.

\section{The controller in closed form}
\label{sec:ctrl}

Let $f_{\mathrm{b}} = 100$~Hz be the design bandwidth and $\omega_{\mathrm{b}} = 2\pi
f_{\mathrm{b}}$. All three channels share one shape, a series-form PID built from a lag-type
integrator, a lead-lag pair and a first-order roll-off:
\begin{equation}
  \Cnorm(s) =
  \underbrace{\frac{s + \omega_{\mathrm{b}}/6}{s}}_{\text{integrator}} \cdot
  \underbrace{\frac{s + \omega_{\mathrm{b}}/3}{s + 3\omega_{\mathrm{b}}}}_{\text{lead-lag}} \cdot
  \underbrace{\frac{10\,\omega_{\mathrm{b}}}{s + 10\,\omega_{\mathrm{b}}}}_{\text{roll-off}}
  = \frac{10\,\omega_{\mathrm{b}}\,(s + \omega_{\mathrm{b}}/6)(s + \omega_{\mathrm{b}}/3)}
         {s\,(s + 3\omega_{\mathrm{b}})(s + 10\,\omega_{\mathrm{b}})} .
  \label{eq:cnorm}
\end{equation}
The plant enters the design through one scalar per channel only. The gain is fixed by requiring
unit open-loop magnitude at the bandwidth,
\begin{equation}
  \kappa_j = \frac{1}{\big|\,G^{\mathrm{c}}_{\mathrm{op},jj}(i\omega_{\mathrm{b}})\,
                       \Cnorm(i\omega_{\mathrm{b}})\,\big|} ,
  \qquad j = 1,2,3 ,
  \label{eq:kappa}
\end{equation}
where $G^{\mathrm{c}}_{\mathrm{op}}$ is the continuous-time design plant of \eqref{eq:designss}, so that
$f_{\mathrm{b}}$ is the $0$~dB crossover of each loop by construction. The applied controller is
the Tustin discretisation of the three scaled shapes,
\begin{equation}
  \Cfb(z) = \diag\big(\kappa_1\Cnorm(s),\ \kappa_2\Cnorm(s),\ \kappa_3\Cnorm(s)\big)
  \Big|_{\,s \,=\, \frac{2}{t_s}\frac{z-1}{z+1}} .
  \label{eq:cfb}
\end{equation}

\medskip\noindent\textbf{Consequence of \eqref{eq:kappa}.}\quad The only $Y_{\mathrm{op}}$
dependence of the controller is through the three numbers $|G^{\mathrm{c}}_{\mathrm{op},jj}
(i\omega_{\mathrm{b}})|$. Between the two design points used by the records this moves $\kappa_1$
by about $5\%$ and $\kappa_3$ by less than $0.1\%$, so the controller is very nearly the same
object at every operating point.

\section{The loop, solved}
\label{sec:loop}

The generator forms the error against a reference $r$ and adds the injected multisine $\fms$ at
the plant input:
\begin{equation}
  e = r - \yh , \qquad
  \ufb = \Cfb\,e , \qquad
  u = \ufb + \fms , \qquad
  \yh = \Gop\,u .
  \label{eq:loopdef}
\end{equation}
Substituting the first three into the fourth and collecting $\yh$,
\begin{equation}
  \yh = \Gop\big(\Cfb (r - \yh) + \fms\big)
  \quad\Longleftrightarrow\quad
  \big(I + \Gop\Cfb\big)\,\yh = \Gop\big(\Cfb\,r + \fms\big) .
  \label{eq:collect}
\end{equation}
Define the single excitation that the loop actually sees, together with the output and input
sensitivities,
\begin{equation}
  w := \Cfb\,r + \fms , \qquad
  \So := \big(I + \Gop\Cfb\big)^{-1} , \qquad
  \Si := \big(I + \Cfb\Gop\big)^{-1} .
  \label{eq:sens}
\end{equation}
Then \eqref{eq:collect} gives the output directly,
\begin{equation}
  \yh = \So\,\Gop\,w = \So\,\Gop\,\Cfb\,r \;+\; \So\,\Gop\,\fms .
  \label{eq:ycl}
\end{equation}
For the applied force, substitute \eqref{eq:ycl} back into \eqref{eq:loopdef},
\begin{equation}
  u = \Cfb\,r - \Cfb\,\So\,\Gop\,w + \fms
    = \big(I - \Cfb\,\So\,\Gop\big)\,w ,
  \label{eq:ustep}
\end{equation}
and apply the push-through identity
\begin{equation}
  I - \Cfb\big(I + \Gop\Cfb\big)^{-1}\Gop = \big(I + \Cfb\Gop\big)^{-1} ,
  \label{eq:pushthrough}
\end{equation}
which follows by left-multiplying both sides by $(I + \Cfb\Gop)$ and cancelling. This yields the
compact form
\begin{equation}
  \boxed{\;u = \Si\,\big(\Cfb\,r + \fms\big) , \qquad
         \yh = \So\,\Gop\,\big(\Cfb\,r + \fms\big) . \;}
  \label{eq:closedform}
\end{equation}
Equation~\eqref{eq:closedform} is the closed-loop form of the system identification course,
Lecture~11, in which the reference enters through the controller and a second excitation enters
at the plant input. The course writes $y = G\So r + \So v$ with a noise term $v$; here $v = 0$,
because the records are generated by a noiseless simulation.

\section{What the loop changes, and what it does not}
\label{sec:conseq}

\subsection{The applied force is the reference shaped by the loop}

By \eqref{eq:closedform} the recorded force is $u = \Si\Cfb\,r + \Si\fms$. The factor $\Cfb$
contains an integrator, so $\Si\Cfb\,r$ carries large low-frequency content whenever the
reference moves. This is the origin of the low-frequency content of the closed-loop records: it
is not disturbance and not an artefact, it is the first term of \eqref{eq:closedform}. The
open-loop records remove it by setting $r = 0$ and $\Cfb = 0$, leaving $u = \fms$.

\subsection{Replaying a recorded force open loop is exact}

The relation $\yh = \Gop u$ in \eqref{eq:loopdef} holds pointwise regardless of what produced
$u$. Hence replaying a recorded $u$ through the same plant without a controller reproduces the
recorded $\yh$ exactly, and no term of \eqref{eq:closedform} is neglected by doing so. The loop
determines \emph{which} $u$ was applied, not whether the replay is valid. Any discrepancy seen in
such a replay is therefore a property of the model used for the replay, not of the missing
controller.

\subsection{Model error inside the loop, exactly}

Let $G$ be the truth and $\hat G$ a candidate model, and let $\Delta := G - \hat G$. Driven open
loop by the same recorded force $u$, the output error is
\begin{equation}
  e_{\mathrm{ol}} = \big(G - \hat G\big)\,u = \Delta\,\Si\,w .
  \label{eq:eol}
\end{equation}
Run instead inside the same loop from the same reference, both outputs take the form
\eqref{eq:ycl} with their own sensitivity, and the difference collapses exactly:
\begin{equation}
  e_{\mathrm{cl}}
  = \big(I + G\Cfb\big)^{-1}G\,w - \big(I + \hat G\Cfb\big)^{-1}\hat G\,w
  = \So\,\Delta\,\hat S_{\mathrm{o}}\,w ,
  \qquad \hat S_{\mathrm{o}} := \big(I + \hat G\Cfb\big)^{-1} ,
  \label{eq:ecl}
\end{equation}
which is an identity, not a linearisation: expand the middle expression over the common
denominators and the cross terms $G\Cfb\hat G$ cancel. Comparing \eqref{eq:eol} and
\eqref{eq:ecl}, both carry one sensitivity factor on the right, and the closed-loop error carries
one \emph{additional} factor $\So$ on the left. The loop therefore rescales the model error by
$\So$; it does not create it and it does not remove it.

\medskip\noindent\textbf{The rescaling is not a suppression in the band of interest.} Evaluating
$\So$ on the frozen design loop gives

\begin{center}
\begin{tabular}{lrrrrrrr}
\toprule
$f$ [Hz] & 1 & 10 & 50 & 100 & 150 & 180 & 500 \\
\midrule
$\sigma_{\max}(\So)$, $Y_{\mathrm{op}} = 0.10$~m & 0.0004 & 0.0214 & 1.0544 & 1.6695 & 1.7983 & 1.8043 & 1.1438 \\
$\sigma_{\max}(\So)$, $Y_{\mathrm{op}} = 0.00$~m & 0.0004 & 0.0217 & 1.0807 & 1.6569 & 1.7937 & 1.8076 & 1.1450 \\
\bottomrule
\end{tabular}
\end{center}

\noindent The loop attenuates strongly only below roughly $45$~Hz. Above that there is a
sensitivity peak of about $1.8$ which spans the whole band $[130,180]$~Hz used for the
augmentation excitation, and in particular the absorber frequency $f_a = 150$~Hz. Closing the
loop around the comparison would therefore \emph{amplify} the discrepancy of interest by roughly
a factor $1.8$, while suppressing the low-frequency part of it by two to three decades. Neither
effect improves the measurement of the absorber, so \eqref{eq:ecl} is an argument for keeping the
comparison open loop, not for restoring the controller.

\section{Noiseless specialisation}
\label{sec:noiseless}

\noindent\textbf{Standing assumption.}\quad The records used throughout this work are generated by
a \emph{noiseless} simulation: there is no measurement noise and no process disturbance, so
$v \equiv 0$ in the notation of the system identification course, Lecture~11. This is a property of
the current data set, not a modelling approximation, and it is deliberate: noise is introduced only
after the noiseless case works. Every result in this section is the course's general result with
$v$ set to zero, and each is shown in both forms so the provenance stays visible. If noise is added
later, this entire section reverts to the general column.

\medskip\noindent\textbf{N1. The predictor is a pure simulation.}\quad The direct method's predictor
(Lecture~11, slide~14) and its noiseless form:
\begin{equation}
  \varepsilon_k = H^{-1}(q)\big[\,y_k - G(q)\,u_k\,\big]
  \qquad\xrightarrow{\ v\,\equiv\,0\ }\qquad
  \varepsilon_k = y_k - G(q)\,u_k .
  \label{eq:n1}
\end{equation}
With $H_{\mathrm{o}} = I$ there is no noise model to estimate, and the requirement that a
Box-Jenkins structure be used to guarantee $\mathcal{S} \in \mathcal{M}$ (slide~19) does not arise.
Equation~\eqref{eq:n1} is the output-error, or simulation-error, objective, which is the objective
the training pipeline already minimises.

\medskip\noindent\textbf{N2. The indirect methods collapse onto the direct method.}\quad The input
decomposition of the two-stage method (slide~11 and slide~30):
\begin{equation}
  u = S_{\mathrm{o}}\,r - K S_{\mathrm{o}}\,v ,
  \qquad u = u_r + u_e
  \qquad\xrightarrow{\ v\,\equiv\,0\ }\qquad
  u_r = u , \quad u_e = 0 .
  \label{eq:n2}
\end{equation}
The component of the input that is driven by the disturbance vanishes, so the reconstruction
$\hat u_r = G_{ur}(q,\hat\theta)\,r$ that the two-stage method performs in its first stage returns
the measured input itself. Stage two then identifies from $(u, y)$, which is the direct method.
The classical indirect construction (slide~29) and the coprime factor construction (slide~32)
reduce in the same way. For the present data these methods are therefore not alternatives to the
direct method; they are the direct method with an additional estimation step that can only add
error.

\medskip\noindent\textbf{N3. The closed-loop bias term vanishes.}\quad The nonparametric estimate
obtained from a spectral analysis (slide~9):
\begin{equation}
  \hat G(e^{i\omega}) =
  \frac{G_{\mathrm{o}}\,\Phi_r - K\,\Phi_v}{\Phi_r + |K|^2\,\Phi_v}
  \qquad\xrightarrow{\ \Phi_v\,\equiv\,0\ }\qquad
  \hat G(e^{i\omega}) = G_{\mathrm{o}} .
  \label{eq:n3}
\end{equation}
The bias that motivates the whole of Lecture~11 is absent, and the degenerate limit
$\hat G \to -1/K$, in which a closed-loop estimate returns the inverse controller instead of the
plant, cannot occur.

\medskip\noindent\textbf{N4. Informativity has a single route, and it is a constraint.}\quad The
input spectrum in closed loop (slide~23):
\begin{equation}
  \Phi_u(\omega) = |S_{\mathrm{o}}|^2\,\Phi_r(\omega) + |K|^2\,\Phi_v(\omega)
  \qquad\xrightarrow{\ \Phi_v\,\equiv\,0\ }\qquad
  \Phi_u(\omega) = |S_{\mathrm{o}}|^2\,\Phi_r(\omega) .
  \label{eq:n4}
\end{equation}
Lecture~11 lists three ways to make closed-loop data informative: a persistently exciting
reference, excitation through the disturbance acting via a controller of sufficient order, or a
time-varying or nonlinear controller. With $v \equiv 0$ the second route is unavailable, and the
controller here is linear and time invariant, so the third is unavailable as well. \emph{All}
identifiability therefore rests on the designed excitation, that is on $r$ and on the injected
multisine $f_{\mathrm{ms}}$. This is the only place where removing the noise makes the problem
harder rather than easier, and it is the formal statement of why the excitation design carries the
whole identifiability burden.

\medskip\noindent\textbf{N5. The loop relations are unchanged.}\quad
Equation~\eqref{eq:closedform} was derived without a disturbance entering the loop, so it is
already the noiseless specialisation and needs no modification.

\medskip\noindent\textbf{Consequence for the identification method.}\quad By \eqref{eq:n1} to
\eqref{eq:n3} the direct method, that is fitting the model to the measured $u$ and $y$ with a
simulation-error objective and no noise model, is exact for this data, and the indirect machinery
of Lecture~11 offers nothing to recover. By \eqref{eq:n4} the burden that remains is experiment
design, not estimation.

\section{Numerical gate}
\label{sec:gate}

Equations~\eqref{eq:cnorm} to \eqref{eq:cfb} were implemented from the formulas alone, with no
MATLAB object and no exported matrices, and applied to the stored $r - \yh$ of two closed-loop
records. The result is compared against the stored $\ufb$ that MATLAB produced. Relative errors
are integrated over $[0,200]$~Hz.

\begin{center}
\begin{tabular}{llrrr}
\toprule
record & channel & $\ufb$ rms [N] & residual rms [N] & relative, $\le 200$~Hz \\
\midrule
\texttt{V1\_standstill\_Yp10} & $X_1$ & 28.63 & $9.01\cdot10^{-7}$ & $4.53\cdot10^{-9}$ \\
\texttt{V1\_standstill\_Yp10} & $X_2$ & 29.17 & $9.03\cdot10^{-7}$ & $4.69\cdot10^{-9}$ \\
\texttt{V1\_standstill\_Yp10} & $Y$   & 24.28 & $6.90\cdot10^{-1}$ & $4.78\cdot10^{-4}$ \\
\texttt{T10\_aprbs\_60}       & $X_1$ & 157.59 & $1.78\cdot10^{-1}$ & $1.27\cdot10^{-4}$ \\
\texttt{T10\_aprbs\_60}       & $X_2$ & 161.03 & $1.83\cdot10^{-1}$ & $1.28\cdot10^{-4}$ \\
\texttt{T10\_aprbs\_60}       & $Y$   & 80.37 & $2.61\cdot10^{-1}$ & $2.05\cdot10^{-4}$ \\
\bottomrule
\end{tabular}
\end{center}

\noindent The spread across rows is set by storage, not by the formulas. The records store
$\ufb$, $\yh$ and $r$ in single precision while MATLAB computed $\ufb$ in double, and the
relative error tracks the float32 quantisation step of the stored position rather than the size
of the signal: the two rows whose stage positions sit at zero, and whose quantisation step is
therefore $9\cdot10^{-13}$~m, agree to $4.5\cdot10^{-9}$, while every row with a step four
decades larger degrades by three to four decades. The residual spectra are flat, which is the
signature of quantisation; a wrong pole, zero, gain or discretisation would place structure at
the frequency it got wrong.

\medskip\noindent\textbf{Conclusion.}\quad Equations~\eqref{eq:cnorm} to \eqref{eq:cfb} are the
controller that generated the records, and \eqref{eq:closedform} may be used in the derivation
as an exact description of how that controller shaped the data.

\end{document}
